\documentclass{article}
\usepackage{amsmath}

\begin{document}

## 1. 인라인 및 디스플레이 수학 환경 (Math Mode)

* mk ➡️ $ $ (Inline Math: 문장 안의 수학 공식)
* dm ➡️ \[ \] (Display Math: 별도 줄에 강조되는 수학 공식)

## 2. 자동 확장 수식 기호 (Autosnippets)
스페이스바나 엔터를 누르지 않아도 입력 즉시 수식으로 변환되는 스니펫들입니다. 속도를 비약적으로 높여줍니다. [5, 6, 7] 

* sr ➡️ ^2 (Square: 제곱, 예: xsr ➡️ x²)
* cb ➡️ ^3 (Cube: 세제곱)
* rd ➡️ ^{} (Power: 임의의 거듭제곱)
* _ ➡️ _{} (Subscript: 아래첨자 자동 중괄호 생성) [8] 

\( a__ \) __ don't work, sometimes, it works. don't know why !!!
\( ;a al ;al :a a_{alb} \)

## 3. 분수 및 심볼 자동 완성 (Fractions & Symbols)
나눗셈과 그리스 문자, 화살표 등 자주 쓰이는 수학 기호 축약어입니다. [9, 10] 

* // ➡️ \frac{}{} (Fraction: 분수 함수 자동 생성)
* lim ➡️ \lim_{n \to \infty} (Limit: 극한 기호)
* inf ➡️ \infty (Infinity: 무한대 기호)
* -> ➡️ \to 또는 \rightarrow (화살표)

lim : don't work.

## 4. 블록 환경 자동 생성 (Environments)
LaTeX의 긴 \begin{} ~ \end{} 블록을 타이핑 없이 생성하는 스니펫입니다. 주로 줄 시작 부분(line_begin)에서 트리거되도록 설정합니다. [7] 

* beg ➡️ \begin{} ... \end{} 환경을 만들고 커서를 중괄호 안에 배치
* ali ➡️ \begin{align*} \end{align*} (수식 정렬 블록)
* enum ➡️ \begin{enumerate} \item \end{enumerate} (번호 매기기 목록)
enum : don't work.

------------------------------
## 💡 가장 중요한 팁: VimTeX 연동 (Math Zone 감지)
Neovim LaTeX 스니펫의 핵심은 "일반 텍스트를 쓸 때"와 "수학 공식 안($ 또는 $$ $$ )에서 쓸 때"를 구별하는 것입니다. [3] 
예를 들어 일반 본문을 적을 때 inf를 입력했다고 \infty로 바뀌면 안 되기 때문입니다. [VimTeX 플러그인](https://github.com/lervag/vimtex)의 vimtex#syntax#in_mathzone() 함수를 LuaSnip 조건문에 연동하면, 오직 수학 모드 안에서만 위의 스니펫들이 자동 작동하도록 세팅할 수 있습니다. [3, 11] 
직접 구현하기 복잡하다면, 깃허브에서 이미 세팅이 완료된 오픈소스 패키지인 [luasnip-latex-snippets.nvim](https://github.com/iurimateus/luasnip-latex-snippets.nvim)을 플러그인 매니저로 설치해 사용하는 것을 추천합니다. [3] 
------------------------------

% -----------------------------------------------------------------
% [테스트 1] 수식 모드 바깥 (Plain Text 구역)
% -----------------------------------------------------------------
% 아래 빈 라인에 'ff'를 타이핑해 보세요.
% 예상 결과: 아무런 자동 트리거도 일어나지 않고, 일반 글자 'ff'만 입력되어야 정상입니다.
% (blink.cmp 후보 창에 friendly-snippets 목록이 뜰 수는 있으나 자동 변환은 안 됨)

여기는 수식 바깥 구역입니다: 


% -----------------------------------------------------------------
% [테스트 2] 인라인 수식 모드 내부 ($ ... $)
% -----------------------------------------------------------------
% 아래 달러 기호($) 사이의 커서 위치(|)에 'ff'를 입력해 보세요.
% 예상 결과: 마지막 'f'를 누르는 순간 탭이나 엔터 없이 바로 $\frac{ }{ }$ 로 변환됩니다.
% 커서가 분모/분자 위치로 들어가면 <C-l> 이나 <Tab>을 눌러 다음 칸으로 점프하는지도 확인하세요.

인라인 수식 테스트: $|$


% -----------------------------------------------------------------
% [테스트 3] 디스플레이 수식 모드 내부 ($$ ... $$)
% -----------------------------------------------------------------
% 아래 디스플레이 수식 모드 한가운데 공간에서 'ff' 또는 '비슷한 수식 단축어'를 입력해 보세요.
% 예상 결과: 인라인 모드와 마찬가지로 수식 감지(Tree-sitter Injection)가 작동하여 자동 변환됩니다.

디스플레이 수식 테스트:
$$
  |
$$

\end{document}
